\section{The equations of the earlier paper}\label{sec:audit}

We first go through the formulas of \cite{tb23}, since the rest of the paper
replaces them. They fall into two groups: kinematic relations between phase,
time and frequency, which are correct, and formulas for an energy loss
$\Delta E$, which are not.

\subsection{Phase, time and frequency}

For a sinusoid of frequency $f$, a delay $\Delta t$ produces the phase
\begin{equation}\label{eq:audit-phase}
\Phi^\circ=360^\circ f\,\Delta t ,
\end{equation}
and \cite{tb23} rewrites this as $\Delta t=\Phi^\circ/(360^\circ f)$ and
$f=\Phi^\circ/(360^\circ\Delta t)$. These are correct, with the
understanding that the last one gives the frequency at which a \emph{given}
delay produces a \emph{given} phase; it does not describe a change of
frequency. The delay that produces one degree is
\[
T(\mathrm{deg})=\frac{1}{360f},
\]
which at $f=5$\,MHz is $555.6$\,ps, as stated in \cite{tb23}. Two of the
phrasings there should be read with care. The equation ``$1^\circ$ phase
shift $=T/360$'' equates an angle with a time; what is meant is that a delay
of $T/360$ produces a phase of $1^\circ$, where $T=1/f$ is the period. The
name ``time distortion'' for $\Delta t$ suggests something more than a delay;
the quantity in \eqref{eq:audit-phase} is just the delay. Finally, the
statement that $T(\mathrm{deg})\propto1/f$ is true of the delay needed for one
degree, and it is the reason why a delay cannot produce the same phase at
two frequencies.

\subsection{The energy formulas}

The paper \cite{tb23} proposes three formulas for an ``infinitesimal loss of
wave energy'',
\begin{equation}\label{eq:audit-dE}
\begin{gathered}
\Delta E=hf\Delta t,\qquad
\Delta E=\frac{2\pi hf_1}{360}\,T(\mathrm{deg}),\\
\Delta E=\frac{2\pi h}{360}\,\frac{T(\mathrm{deg})}{\Delta t},
\end{gathered}
\end{equation}
with $h$ Planck's constant, and evaluates the first one for $f=4{,}986{,}111.11$\,Hz and
$\Delta t=555$\,ps. Four things go wrong.

\emph{Units.} Planck's constant is an action, measured in J\,s; $hf$ is an
energy and $hf\Delta t$ is again an action. The same holds for the other two
formulas in \eqref{eq:audit-dE}, in which $T(\mathrm{deg})/\Delta t$ is a pure
number and $f_1T(\mathrm{deg})$ is the number $1/360$. None of the three
expressions is an energy.\footnote{If $T(\mathrm{deg})$ is read as the angle
$1^\circ=\pi/180$ rather than as a time, the second formula has units of
energy, but it then reads $\Delta E=(\pi/180)\,hf_1\,(2\pi/360)$, a fixed
multiple of $hf_1$ without any connection to the delay.}

\emph{Arithmetic.} The product evaluated in \cite{tb23} is
\[
(6.626\times10^{-34})(4.986\times10^{6})(555\times10^{-12})
=1.834\times10^{-36},
\]
in J\,s, not $1.848\times10^{-27}$\,J. The stated value is larger by a factor
of about $10^{9}$, and it would amount to $56\%$ of the energy
$E_1=hf_1=3.313\times10^{-27}$\,J of a quantum at $5$\,MHz, which can hardly
be called infinitesimal.

\emph{A fixed phase shift does not change the frequency.} The second example
of \cite{tb23} assigns to a $1^\circ$ phase shift the new frequency
$f_2=f_1-f_1/360=4{,}986{,}111.11$\,Hz. Multiplying a wave by a constant
$e^{i\theta}$, or delaying it, leaves its frequency unchanged
(\cref{prop:static}). A frequency change of $f_1/360$ occurs when the phase
\emph{slips} by one degree in every period, that is, when the phase grows
at the rate $\dot\theta=2\pi f_1/360$ (\cref{prop:slip}). In that case the
energy of a quantum changes by
\[
E_1-E_2=\frac{hf_1}{360}=9.20\times10^{-30}\ \mathrm J ,
\]
and this is the only energy difference of the size suggested by the example.
The paper also gives $E_2=hf_2\approx3.313\times10^{-27}$\,J; the correct
value is $3.304\times10^{-27}$\,J, and the difference from $E_1$ is the number
just computed, not $1.848\times10^{-27}$\,J.

\emph{What $hf$ measures.} The energy $hf$ is that of one quantum. The energy
of a classical wave is fixed by its amplitude, which the formulas
\eqref{eq:audit-dE} do not contain. A loss caused by a passive element is a
fraction of the incident energy, whether the wave is classical or quantised,
and the quantity to compute is therefore the transmission $\abs{H}^2$, or,
per quantum, the probability of absorption.

\subsection{What remains}

\Cref{tab:audit} lists the corrections. The kinematic part of \cite{tb23}
stands. The energy part asks a good question in the wrong variables: the
phase shift that has a price is the one applied uniformly over a band, the
price is a fraction of the energy rather than a multiple of $h$, and its
size is governed by the bandwidth. Sections~\ref{sec:passive}--\ref{sec:waves}
compute it.

\begin{table}[t]
\caption{The formulas of \cite{tb23} and their status.}
\label{tab:audit}
\centering
\small
\begin{tabular}{@{}p{0.43\columnwidth}p{0.51\columnwidth}@{}}
\toprule
Formula in \cite{tb23} & Status \tabularnewline
\midrule
\raggedright $\Phi^\circ=360^\circ f\Delta t$ and its inversions & \raggedright correct\tabularnewline
\raggedright $T(\mathrm{deg})=1/(360f)$, $555$\,ps at $5$\,MHz & \raggedright correct\tabularnewline
\raggedright ``$1^\circ$ phase shift $=T/360$'' & \raggedright angle equated with time; a delay $T/360$ gives $1^\circ$\tabularnewline
\raggedright $\Delta E=hf\Delta t$ & \raggedright an action, in J\,s; its value is $1.83\times10^{-36}$, not $1.848\times10^{-27}$\tabularnewline
\raggedright $\Delta E=(2\pi hf_1/360)T(\mathrm{deg})$ & \raggedright an action; equals $2\pi h/360^2$\tabularnewline
\raggedright $\Delta E=(2\pi h/360)\,T(\mathrm{deg})/\Delta t$ & \raggedright an action\tabularnewline
\raggedright $f_2=f_1-f_1/360$ after a $1^\circ$ shift & \raggedright only if the phase slips $1^\circ$ per period\tabularnewline
\raggedright $E_2\approx3.313\times10^{-27}$\,J & \raggedright $3.304\times10^{-27}$\,J; $E_1-E_2=9.20\times10^{-30}$\,J\tabularnewline
\bottomrule
\end{tabular}
\end{table}
